\section{Threat Model}
\label{sec:threat-model}

\draftnote{DRAFT, but the underlying design/table is final and
implemented -- see \texttt{puffin/attacks/README.md}.}

We distinguish mechanisms by what they actually release to an
adversary, and run only the attacks that match that release -- an
attack that assumes an artifact the mechanism does not release
overstates risk and is not meaningful evidence against it.

\begin{table}[h]
\centering
\caption{Per-mechanism threat model: released artifact and the attacks
that apply to it.}
\label{tab:threat-model}
\small
\begin{tabular}{@{}lll@{}}
\toprule
Mechanism & Released artifact & Applicable attacks \\
\midrule
Mo-PAE & encoder features $f$ & external re-id on $f$; MIA on the \\
       & (worst case: the encoder itself) & utility model; decoder reconstruction \\
TrajGAN & synthetic trajectories & trajectory-user linking (TUL); MIA \\
GI-DP & perturbed location sequences & TUL; re-id directly on sequences \\
\bottomrule
\end{tabular}
\end{table}

For Mo-PAE, every re-identification risk number we report is from a
\emph{fresh external attacker}, trained strictly after the mechanism
is frozen, never the in-model adversary the mechanism's own encoder was
trained to fool -- reporting the in-model adversary's own residual
accuracy as a privacy measure would overstate protection, since the
encoder was optimized specifically against that one adversary and
nothing else.

\subsection{The naive (primary) external attacker}
\label{sec:threat-model-naive}

Our primary evaluation standard, used throughout this paper except
where explicitly marked otherwise, is a \emph{hardened, group-agnostic}
external attacker: for each of five random seeds, we train three
candidate architectures -- an MLP on the flattened window, a GRU over
the window's sequence, and a gradient-boosted tree ensemble on the
flattened window's raw features -- and take, per user, the
element-wise maximum top-1/top-5 accuracy across all three
architectures, then average that per-user maximum across the five
seeds. This reports, for each user, the risk posed by the \emph{best
attacker we tried for that specific user}, not one globally-chosen
architecture, following the design in \texttt{puffin/attacks/reid.py}.

\subsection{The retired adaptive adversary}
\label{sec:threat-model-adaptive}

An earlier phase of this project also studied an \emph{adaptive}
adversary explicitly aware of (a) a mechanism's fairness objective and
(b) each user's true behavioral group (entropy decile), on the
hypothesis that such an adversary would recover more risk for the
group a mechanism protects most than the naive attacker above reports.
We retired this adversary as a primary evaluation standard after
finding (Section~\ref{sec:adaptive-appendix}) that its apparent
advantage is mostly explained by a smaller candidate pool and oracle
group knowledge rather than genuine attacker sophistication. We retain
it here only as a robustness appendix.
